% ============================================================================
\section{Problème et objectifs}
% ============================================================================

\subsection{La question du projet et notre scénario}

Le projet consiste à construire une petite plateforme d'exécution distribuée et
à déterminer, pour un lot de tâches donné, \textbf{sur quelle machine virtuelle
exécuter chaque tâche}.

L'énoncé sépare quatre éléments, que l'on retrouve dans notre architecture
(section~\ref{sec:architecture}) : les \textbf{workloads} (programmes compilés),
l'\textbf{infrastructure} (VM Google Compute Engine), le \textbf{scheduler}
(qui \emph{décide} où va chaque tâche) et l'\textbf{orchestrateur} (qui
\emph{exécute} réellement les tâches et collecte les résultats).

Pour donner un sens à l'optimisation, nous nous plaçons dans la situation
(fictive) d'une \textbf{petite entreprise de traitement par lots}, qui exécute
des calculs pour ses clients : calcul numérique, traitement mémoire, calcul
matriciel, tri de données. Ce n'est pas un service web : un résultat livré
quelques secondes plus tard est acceptable, mais chaque euro dépensé inutilement
se répète à chaque lot. Ses priorités, par ordre d'importance :
\begin{enumerate}
  \item \textbf{réduire la facture cloud} ;
  \item \textbf{conserver un délai raisonnable} ;
  \item pouvoir traiter des lots de quelques tâches à plusieurs milliers ;
  \item rester simple à maintenir et à expliquer.
\end{enumerate}

\subsection{Objectif et contrainte}

Nous évaluons une allocation par deux grandeurs :
\begin{itemize}
  \item $D$, le \textbf{makespan} : le temps entre le début du lot et la fin de la
        dernière tâche ;
  \item $C$, le \textbf{coût} : pour chaque VM utilisée, son temps d'utilisation
        (avec la facturation minimale de 60~s de GCP) multiplié par son prix horaire.
\end{itemize}

Ces deux grandeurs n'ont pas la même unité. Nous les normalisons par celles de
\textbf{Round Robin} (RR), notre référence, et les combinons en un score :
\begin{equation}
  S = 0{,}6 \cdot \frac{C}{\Crr} + 0{,}4 \cdot \frac{D}{\Drr}
  \qquad \text{(plus bas = mieux ; RR vaut 1)}
  \label{eq:score}
\end{equation}

Le poids de 60\,\% sur le coût traduit la priorité de l'entreprise. Ce n'est pas
une valeur imposée : c'est un \textbf{paramètre} du scheduler (option
\code{-wcost}). Pour éviter qu'une solution très économique mais très lente soit
retenue, nous ajoutons une \textbf{contrainte de délai} :
\begin{equation}
  D \le 1{,}25 \cdot \Drr
  \label{eq:deadline}
\end{equation}
Une allocation qui la viole est dite \emph{non admissible}, quel que soit son
score. Le seuil de 125\,\% est un choix de conception, lui aussi paramétrable
(\code{-deadline}).

\subsection{Ce qui a changé depuis notre première version}

Une première version du projet proposait un algorithme « V8 » (glouton pondéré
puis recherche locale), évalué uniquement en simulation, et dont le modèle de
durée sous-estimait les temps réels d'un facteur 3. Nous avons repris le projet
en partant des mesures. Le scénario, le score $S$ et la contrainte de délai sont
conservés ; le reste a été refait :

\begin{table}[H]
\centering\small
\begin{tabularx}{\textwidth}{>{\raggedright\arraybackslash}p{3.3cm}XX}
\toprule
 & \textbf{Première version} & \textbf{Version actuelle} \\
\midrule
Mesures & 180 mesures, VM3 recopiée sur VM2 & 4 VM mesurées séparément + mesure du parallélisme réel \\
Modèle de durée & médiane du benchmark (erreur $\times 3$) & médiane $\times$ contention + surcoût SSH (erreur $\approx \pm 10\,\%$) \\
Algorithme & V8 : recherche locale qui minimise le makespan & \sel{} : choix exact des VM à allumer + Max-Min, minimise $S$ \\
Exécution distribuée & non démontrée & orchestrateur Go, des milliers de tâches exécutées \\
Résultats & simulés uniquement & simulés \emph{et} mesurés sur GCP \\
\bottomrule
\end{tabularx}
\caption{Évolution du projet.}
\end{table}
